\section*{第 \ref{ch:g6:latitude-longitude} 章——纬度、经度与时区}

\begin{solution}{exo:g6:latitude-longitude:1}
“\omterm{def:g6:latitude-longitude:latitude}{纬度}从\omterm{def:g3:continents-and-oceans:equator}{赤道}开始计量，范围为 $0^\circ$ 至
$90^\circ$（北纬或南纬）；\omterm{def:g6:latitude-longitude:longitude}{经度}从\omterm{def:g6:latitude-longitude:longitude}{本初子午线}开始计量，范围为 $0^\circ$ 至
$180^\circ$（东经或西经）。”
\end{solution}

\begin{solution}{exo:g6:latitude-longitude:2}
\omterm{def:g3:continents-and-oceans:poles}{北极}：$90^\circ$\,N。\omterm{def:g3:continents-and-oceans:poles}{南极}：$90^\circ$\,S。\omterm{def:g3:continents-and-oceans:equator}{赤道}上的
一个城镇：$0^\circ$。
\end{solution}

\begin{solution}{exo:g6:latitude-longitude:3}
悉尼：\omterm{def:g3:continents-and-oceans:equator}{南半球}，东经。纽约：\omterm{def:g3:continents-and-oceans:equator}{北半球}，西经。
开罗：\omterm{def:g3:continents-and-oceans:equator}{北半球}，东经。里约热内卢：\omterm{def:g3:continents-and-oceans:equator}{南半球}，
西经。
\end{solution}

\begin{solution}{exo:g6:latitude-longitude:4}
(a) 对。(b) 错：\omterm{def:g6:latitude-longitude:latitude}{纬线}越靠近两极越短，最长的是
\omterm{def:g3:continents-and-oceans:equator}{赤道}。(c) 对。(d) 错：它经过格林尼治，
位于伦敦。
\end{solution}

\begin{solution}{exo:g6:latitude-longitude:5}
开罗。纽约。东京位于约 $36^\circ$\,N、$140^\circ$\,E（略
高于\omterm{def:g6:latitude-longitude:latitude}{纬线} $30^\circ$\,N，略超过\omterm{def:g6:latitude-longitude:longitude}{经线} $120^\circ$\,E，
未达 $150^\circ$\,E）。
\end{solution}

\begin{solution}{exo:g6:latitude-longitude:6}
$45 \div 15 = 3$：两地正午时间相差 3 小时。
先出现正午的是城镇~B，它更靠东。
\end{solution}

\begin{solution}{exo:g6:latitude-longitude:7}
$73 \div 15 \approx 4.9$，四舍五入为 5：该\omterm{def:g6:latitude-longitude:time-zone}{时区}为 $+5$。
格林尼治正午时，孟买为 $12 + 5 = 17$，即下午 5 点。
\end{solution}

\begin{solution}{exo:g6:latitude-longitude:8}
$30 \times 111 = 3\,330$~千米。\omterm{def:g3:continents-and-oceans:poles}{北极}位于 $90^\circ$：
$90 \times 111 = 9\,990$~千米，约 10\,000~千米。
\end{solution}

\begin{solution}{exo:g6:latitude-longitude:9}
两地相差 $9 - (-5) = 14$ 小时。$8 + 14 = 22$：
此时东京为周一晚上 10 点。
\end{solution}

\begin{solution}{exo:g6:latitude-longitude:10}
她说得对。地球自西向东自转，因此东边的城镇先被
带向太阳：它们先看到太阳升起、到达最高点，
然后落山，都比更西边的城镇早。
\end{solution}

\begin{solution}{exo:g6:latitude-longitude:11}
(a) $42 \times 3 = 126$ 秒，略多于 2 分钟。(b)
$126 \div 240 \approx 0.5$：约半度\omterm{def:g6:latitude-longitude:longitude}{经度}。(c)
$0.5 \times 111 \approx 55$~千米（取 $0.525$ 时约 58~千米）。
误差约五十千米，足以错过一座岛，或一头撞上
本以为远在别处的\omterm{def:g1:landscapes-around-us:water}{海岸}。
\end{solution}

\begin{solution}{exo:g6:latitude-longitude:12}
(a) $18 - 9 = 9$：周一格林尼治时间上午 9 点。(b) $9 + 10 = 19$：
周一格林尼治时间下午 7 点。(c) $19 - 8 = 11$：周一
洛杉矶时间上午 11 点。飞机以东京时间周一下午 6 点起飞，
并以洛杉矶时间同一个周一上午 11 点降落：到达时的钟面
显示的时间比出发时更早，因为飞机向东越过了\omterm{def:g6:latitude-longitude:date-line}{国际日期变更线}，
于是日期倒退了一天。
\end{solution}

\begin{solution}{pb:g6:latitude-longitude:1}
\textbf{1.} 安娜（里约热内卢，$23^\circ$\,S）与米娅（奥克兰，
$37^\circ$\,S）。

\textbf{2.} 安娜与萨姆（$43^\circ$\,W 和 $118^\circ$\,W）。
莉娜位于\omterm{def:g6:latitude-longitude:longitude}{本初子午线}上。

\textbf{3.} 莉娜：伦敦位于 $52^\circ$\,N，比其他所有城市
都更远离\omterm{def:g3:continents-and-oceans:equator}{赤道}。

\textbf{4.} 米娅：奥克兰位于 $175^\circ$\,E，仅差 $180 - 175 = 5$ 度
便到 $180^\circ$ \omterm{def:g6:latitude-longitude:longitude}{经线}。

\textbf{5.} $360^\circ \div 24 = 15^\circ$。

\textbf{6.} 伦敦：$0$。里约热内卢：$43 \div 15 \approx 2.9$，西经，
故为 $-3$。东京：$140 \div 15 \approx 9.3$，东经，故为 $+9$。奥克兰：
$175 \div 15 \approx 11.7$，东经，故为 $+12$。洛杉矶：
$118 \div 15 \approx 7.9$，西经，故为 $-8$。

\textbf{7.} 里约热内卢：$20 - 3 = 17$，周五下午 5 点。东京：
$20 + 9 = 29$；$29 - 24 = 5$：周六上午 5 点。

\textbf{8.} 东京位于很靠东的位置，其时钟
比伦敦快 9 小时。伦敦晚上 8 点时，东京已过午夜，
所以东京的周六已经开始。

\textbf{9.} $20 - 8 = 12$：周五中午 12 点。

\textbf{10.} 奥克兰比洛杉矶快 $12 - (-8) = 20$ 小时。

\textbf{11.} 船向东越过\omterm{def:g6:latitude-longitude:date-line}{国际日期变更线}时，
日期倒退回前一天。

\textbf{12.} $20 + 12 = 32$；$32 - 24 = 8$：通话开始时，
奥克兰时间为\textbf{周六上午 8 点}。米娅在早餐时间加入通话，
而萨姆正在周五吃午饭。
\end{solution}
